\documentclass[../DoAn.tex]{subfiles}
\begin{document}

\appendixtitleon
\appendixtitletocon

\section{Hướng dẫn cài đặt môi trường phát triển}
\label{sec:appendix-a1}

Phần này cung cấp hướng dẫn chi tiết từng bước để thiết lập môi trường phát triển cho hệ thống SMART GRADING trên máy tính cá nhân.

\subsection{Cài đặt công cụ phát triển backend}
\label{sec:appendix-a1-1}

Để phát triển backend server, cần cài đặt các công cụ sau. Đầu tiên là cài đặt Node.js phiên bản 20 LTS bằng cách tải installer từ trang chủ nodejs.org và chạy installer, sau khi cài đặt xong kiểm tra bằng lệnh node --version trong terminal. Thứ hai là cài đặt MongoDB phiên bản 7.0 bằng cách tham khảo hướng dẫn tại trang chủ mongodb.com cho hệ điều hành tương ứng, có thể sử dụng MongoDB Community Edition miễn phí hoặc MongoDB Atlas (cloud) cho môi trường phát triển. Thứ ba là cài đặt Git bằng cách tải và cài đặt từ trang git-scm.com, Git được sử dụng để quản lý mã nguồn và làm việc với repository.

Sau khi cài đặt các công cụ cơ bản, tiến hành cài đặt project backend bằng các bước sau. Bước 1 là mở terminal và navigate đến thư mục muốn chứa project. Bước 2 là clone repository bằng lệnh git clone <repository-url>. Bước 3 là navigate vào thư mục server bằng lệnh cd server. Bước 4 là chạy npm install để cài đặt tất cả dependencies. Bước 5 là tạo file .env trong thư mục server với nội dung mẫu được cung cấp trong file .env.example. Bước 6 là chạy npm run dev để khởi động server ở chế độ phát triển.

\subsection{Cài đặt môi trường phát triển ứng dụng mobile}
\label{sec:appendix-a1-2}

Để phát triển ứng dụng Flutter, cần cài đặt các công cụ sau. Đầu tiên là cài đặt Flutter SDK bằng cách tải Flutter SDK từ flutter.dev, giải nén vào thư mục mong muốn (ví dụ: C:\\src\\flutter), và thêm đường dẫn flutter/bin vào biến môi trường PATH. Thứ hai là cài đặt Android Studio bao gồm Android SDK, Android SDK Platform-Tools, và Android SDK Build-Tools. Thứ ba là cài đặt Xcode (chỉ dành cho macOS) để phát triển ứng dụng iOS.

Sau khi cài đặt, tiến hành cài đặt project mobile bằng các bước sau. Bước 1 là mở terminal và navigate đến thư mục chứa project. Bước 2 là chạy flutter pub get để cài đặt tất cả dependencies của project. Bước 3 là kết nối thiết bị Android hoặc iOS hoặc mở emulator. Bước 4 là chạy flutter run để khởi chạy ứng dụng trên thiết bị.

\subsection{Cài đặt môi trường phát triển ứng dụng web}
\label{sec:appendix-a1-3}

Để phát triển ứng dụng React, cần cài đặt Node.js phiên bản 20 LTS (đã cài ở bước trước) và npm hoặc yarn là package manager đi kèm Node.js. Tiến hành cài đặt project web bằng các bước sau. Bước 1 là navigate đến thư mục chứa project. Bước 2 là chạy npm install để cài đặt tất cả dependencies. Bước 3 là tạo file .env với biến môi trường VITE\_API\_BASE\_URL trỏ đến backend server. Bước 4 là chạy npm run dev để khởi động development server.

\section{Cấu trúc thư mục project}
\label{sec:appendix-a2}

Phần này mô tả cấu trúc thư mục của các thành phần trong hệ thống SMART GRADING.

Cấu trúc thư mục backend bao gồm các thư mục chính sau. Thư mục src chứa mã nguồn chính, trong đó src/config chứa các file cấu hình như database.js và env.js, src/controllers chứa các route controllers, src/middlewares chứa các Express middlewares, src/models chứa các Mongoose models, src/routes chứa các định nghĩa route, src/services chứa các business logic services, src/utils chứa các utility functions, src/validations chứa các Joi validation schemas, src/app.js là Express app configuration, và src/index.js là entry point. Thư mục tests chứa các unit test và integration test. File .env chứa các biến môi trường.

Cấu trúc thư mục mobile bao gồm các thư mục chính sau. Thư mục lib chứa mã nguồn chính, trong đó lib/core chứa constants, config, network và utilities, lib/domain chứa entities, models, OMR engine và repositories, lib/presentation chứa blocs, pages, widgets và screens, lib/di chứa các dependency injection setup. File pubspec.yaml chứa danh sách dependencies. Thư mục test chứa các widget test và bloc test.

Cấu trúc thư mục web bao gồm các thư mục chính sau. Thư mục src chứa mã nguồn chính, trong đó src/components chứa các React components, src/screens chứa các page components, src/hooks chứa các custom hooks, src/services chứa các API services, src/stores chứa các Zustand stores, src/types chứa các TypeScript type definitions. File package.json chứa danh sách dependencies. Thư mục src/\_\_tests\_\_ chứa các test files.

\section{Danh sách các API endpoints}
\label{sec:appendix-a3}

Phần này cung cấp danh sách các API endpoints chính của hệ thống SMART GRADING.

Nhóm Authentication bao gồm POST /api/auth/register để đăng ký tài khoản mới, POST /api/auth/login để đăng nhập và nhận JWT token, POST /api/auth/refresh để làm mới access token, và POST /api/auth/logout để đăng xuất.

Nhóm Users bao gồm GET /api/users để lấy danh sách người dùng (admin only), GET /api/users/:id để lấy thông tin một người dùng, PUT /api/users/:id để cập nhật thông tin người dùng, và DELETE /api/users/:id để xóa người dùng (admin only).

Nhóm Exams bao gồm GET /api/exams để lấy danh sách đề thi, POST /api/exams để tạo đề thi mới, GET /api/exams/:id để lấy chi tiết một đề thi, PUT /api/exams/:id để cập nhật đề thi, DELETE /api/exams/:id để xóa đề thi, POST /api/exams/:id/versions để tạo các mã đề từ file LaTeX, và GET /api/exams/:id/versions để lấy danh sách mã đề.

Nhóm Submissions bao gồm POST /api/submissions để nộp bài thi, GET /api/submissions để lấy danh sách kết quả, GET /api/submissions/:id để lấy chi tiết một kết quả, và POST /api/submissions/:id/appeal để nộp đơn phúc tra.

Nhóm AI bao gồm POST /api/ai/generate-questions để sinh câu hỏi tự động, POST /api/ai/analyze-results để phân tích kết quả thi, và POST /api/ai/chat để giao tiếp với chatbot AI.

\section{Kịch bản kiểm thử chi tiết}
\label{sec:appendix-a4}

Phần này cung cấp các kịch bản kiểm thử chi tiết cho các chức năng quan trọng của hệ thống.

Kịch bản kiểm thử chức năng đăng nhập bao gồm ba trường hợp. Trường hợp thành công là nhập email và password hợp lệ, hệ thống đăng nhập thành công và chuyển hướng đến trang chủ. Trường hợp thất bại với sai password là nhập email đúng nhưng password sai, hệ thống hiển thị thông báo lỗi và không cho đăng nhập. Trường hợp thất bại với tài khoản không tồn tại là nhập email chưa đăng ký, hệ thống hiển thị thông báo tài khoản không tồn tại.

Kịch bản kiểm thử chức năng quét OMR bao gồm ba trường hợp. Trường hợp thành công là chụp ảnh phiếu rõ ràng trong điều kiện ánh sáng tốt, hệ thống nhận diện đúng và hiển thị kết quả chấm điểm. Trường hợp ảnh mờ là chụp ảnh phiếu bị mờ hoặc thiếu sáng, hệ thống hiển thị cảnh báo và yêu cầu chụp lại. Trường hợp tô đúp là thí sinh tô nhiều hơn một đáp án cho một câu, hệ thống hiển thị cảnh báo double-fill và yêu cầu giáo viên xác nhận.

Kịch bản kiểm thử chức năng tạo đề thi bao gồm hai trường hợp. Trường hợp thành công là upload file LaTeX hợp lệ và hệ thống tạo thành công các mã đề, hệ thống hiển thị danh sách mã đề và cho phép tải về PDF. Trường hợp thất bại với file LaTeX không hợp lệ là upload file có lỗi cú pháp LaTeX, hệ thống hiển thị thông báo lỗi từ AMC và không tạo được mã đề.

\end{document}
